% =====================================================================
%  Module Data Lakehouse -- TP 2 : Données semi-structurées et texte
%  Compilation : pdflatex (2 passes)
% =====================================================================
\documentclass[11pt,a4paper]{article}

\usepackage[utf8]{inputenc}
\usepackage[T1]{fontenc}
\usepackage[french]{babel}
\usepackage{lmodern}
\usepackage[babel=true]{microtype}

\usepackage[a4paper,margin=2.3cm,headheight=14pt]{geometry}
\usepackage{fancyhdr}
\usepackage{titlesec}
\usepackage{needspace}
\usepackage{enumitem}
\usepackage{amssymb}

\usepackage[table]{xcolor}
\usepackage{booktabs}
\usepackage{tabularx}
\usepackage{array}
\usepackage{tikz}
\usetikzlibrary{positioning,arrows.meta,shapes.geometric,babel}
\usepackage{listings}

\usepackage[htt]{hyphenat}
\usepackage[hidelinks]{hyperref}
\usepackage{xurl}
\usepackage{caption}
\usepackage{float}
\usepackage{graphicx}
% Captures d'écran : <nom>.png à la racine du projet (Overleaf) ou dans images/tp2/ ;
% tant que le fichier manque, un cadre indique le nom attendu.
\newcommand{\capture}[3]{%
  \begin{figure}[H]\centering
  \IfFileExists{#1.png}%
    {\includegraphics[width=0.92\textwidth,height=0.42\textheight,keepaspectratio]{#1.png}}%
    {\IfFileExists{images/tp2/#1.png}%
      {\includegraphics[width=0.92\textwidth,height=0.42\textheight,keepaspectratio]{images/tp2/#1.png}}%
      {\fbox{\parbox[c][3.2cm][c]{0.88\textwidth}{\centering\footnotesize
         Capture à ajouter : \texttt{\detokenize{#1}.png}}}}}%
  \caption{#2}\label{#3}
  \end{figure}}
\captionsetup{font=small,labelfont=bf,labelsep=period}

\setlength{\parskip}{0.5em}
\setlength{\parindent}{0pt}
\setlength{\emergencystretch}{3em}

% ---------- Couleurs ----------
\definecolor{lhblue}{HTML}{1F4E79}
\definecolor{lhteal}{HTML}{0F766E}
\definecolor{lhred}{HTML}{B91C1C}
\definecolor{lhviolet}{HTML}{6D28D9}
\definecolor{lhgray}{HTML}{F3F4F6}
\definecolor{lhbronze}{HTML}{A16207}
\definecolor{lhsilver}{HTML}{64748B}
\definecolor{lhgold}{HTML}{CA8A04}
\definecolor{lhfonce}{HTML}{17394D}
\definecolor{lhvert}{HTML}{5FB89D}
% Logos des logiciels (docs/images/logos/, voir LISEZMOI.md) ; sans le fichier, rien ne s'affiche
\newcommand{\logo}[2][1.15em]{\IfFileExists{images/logos/#2.png}{\raisebox{-0.22em}{\includegraphics[height=#1]{images/logos/#2.png}}\hspace{0.3em}}%
  {\IfFileExists{logos/#2.png}{\raisebox{-0.22em}{\includegraphics[height=#1]{logos/#2.png}}\hspace{0.3em}}{}}}
% Pictogrammes de la figure du pipeline (mêmes dessins que dans le Support 0)
\tikzset{
  pics/fichier/.style={code={%
    \fill[white,draw=lhfonce,line width=0.7pt] (-0.32,-0.42) -- (0.18,-0.42) -- (0.32,-0.28) -- (0.32,0.42) -- (-0.32,0.42) -- cycle;
    \draw[lhfonce,line width=0.7pt] (0.18,-0.42) -- (0.18,-0.28) -- (0.32,-0.28);
    \node[font=\sffamily\bfseries\tiny,text=lhfonce] at (0,0.05) {#1};}},
  pics/table/.style={code={%
    \node[cylinder,shape border rotate=90,aspect=0.28,draw=#1!75!black,fill=#1!35,
          minimum width=0.85cm,minimum height=0.95cm,line width=0.6pt] {};}},
  pastille/.style={circle,fill=lhvert,draw=white,line width=1pt,text=white,
                   font=\sffamily\bfseries\scriptsize,inner sep=1.6pt},
}


% ---------- Titres, en-têtes ----------
\titleformat{\section}{\needspace{6\baselineskip}\Large\bfseries\color{lhblue}}{\thesection.}{0.6em}{}
\titleformat{\subsection}{\large\bfseries\color{lhblue}}{\thesubsection}{0.6em}{}
\pagestyle{fancy}
\fancyhf{}
\fancyhead[L]{\small\color{lhblue}Module Data Lakehouse}
\fancyhead[R]{\small\color{lhblue}TP 2 -- Données semi-structurées et texte}
\fancyfoot[C]{\small\thepage}

% ---------- Encadrés : texte simple, sans boîte ----------
% Un titre en gras introduit le paragraphe ; seul « Attention » est en rouge.
\newenvironment{resultat}{\par\medskip\noindent\textbf{Résultat attendu :}\ }{\par\medskip}
\newenvironment{astuce}{\par\medskip\noindent\textbf{Astuce :}\ }{\par\medskip}
\newenvironment{attention}{\par\medskip\noindent\textbf{\textcolor{lhred}{Attention :}}\ }{\par\medskip}

% ---------- Code ----------
\lstset{
  basicstyle=\ttfamily\small, columns=fullflexible, keepspaces=true,
  breaklines=true, showstringspaces=false, upquote=true,
  frame=single, framesep=5pt, xleftmargin=6pt, xrightmargin=6pt,
  backgroundcolor=\color{lhgray}, rulecolor=\color{black!15},
  commentstyle=\color{black!55}\itshape,
  keywordstyle=\color{lhblue}\bfseries, stringstyle=\color{black!70},
  literate=
   {é}{{\'e}}1 {è}{{\`e}}1 {ê}{{\^e}}1 {à}{{\`a}}1 {ç}{{\c c}}1
   {ù}{{\`u}}1 {î}{{\^i}}1 {ô}{{\^o}}1 {É}{{\'E}}1 {Ê}{{\^E}}1
}
\lstdefinestyle{shell}{language=bash,backgroundcolor=\color{black!85},
  basicstyle=\ttfamily\small\color{white},rulecolor=\color{black!85},
  commentstyle=\color{black!35}\itshape,keywordstyle=\color{white}}

\newcommand{\cmd}[1]{\texttt{#1}}
\newcommand{\cocher}{$\square$\ }
\newcommand{\etape}[1]{\tikz[baseline=(n.base)]\node[circle,fill=lhvert,text=white,
  font=\sffamily\bfseries\scriptsize,inner sep=1.4pt](n){#1};}
\pdfstringdefDisableCommands{\def\cmd#1{#1}}
\hypersetup{pdftitle={TP 2 -- Texte et JSON},pdfauthor={Anouar Riadsolh}}

\newenvironment{notion}{\par\medskip\noindent\textbf{\textcolor{lhblue}{Notion clé :}}\ }{\par\medskip}
\newcommand{\nb}[1]{\textcolor{lhblue}{\texttt{#1}}}

\lstdefinelanguage{json}{morestring=[b]", stringstyle=\color{black!75}, literate=
   {é}{{\'e}}1 {è}{{\`e}}1 {à}{{\`a}}1}

% =====================================================================
\begin{document}
% =====================================================================

\begin{tikzpicture}[remember picture,overlay]
  \fill[lhblue] (current page.north west) rectangle ([yshift=-5.2cm]current page.north east);
  \node[anchor=north west,text=white,font=\sffamily] at ([xshift=2.3cm,yshift=-1.3cm]current page.north west)
    {Module Data Lakehouse};
  \node[anchor=north west,text=white,font=\sffamily\bfseries\fontsize{26}{30}\selectfont] at ([xshift=2.3cm,yshift=-2.2cm]current page.north west)
    {TP 2 -- Texte et JSON};
  \node[anchor=north west,text=white!85!lhblue,font=\sffamily\large] at ([xshift=2.3cm,yshift=-3.6cm]current page.north west)
    {Avis clients, suivi des colis, modélisation dimensionnelle avec dbt et analyse du texte};
  \node[anchor=north east,text=white!85!lhblue,font=\sffamily] at ([xshift=-2.3cm,yshift=-4.25cm]current page.north east)
    {Pr.~Anouar Riadsolh};
\end{tikzpicture}
\thispagestyle{empty}
\vspace*{4.2cm}

{\large\bfseries\color{lhblue}En bref}\par\smallskip
\begin{tabularx}{\linewidth}{@{}>{\bfseries}l X@{}}
\toprule
Objectif & Faire entrer dans le Lakehouse du texte libre et des messages JSON, construire une couche
           Gold en schéma en constellation avec dbt, puis repérer automatiquement les clients mécontents. \\
Durée    & 2 séances de 3 h (+ environ 2 h de travail personnel). \\
Prérequis & TP 1 terminé, avec les tables Silver \cmd{orders}, \cmd{order\_items}, \cmd{customers},
            \cmd{products} et \cmd{sellers}. \\
À rendre & Un compte rendu, les notebooks exécutés, le projet dbt et le tableau de bord
           (section~\ref{sec:livrables}). \\
\bottomrule
\end{tabularx}

% =====================================================================
\section{Le scénario}
% =====================================================================

Le TP 1 a montré qu'Olist livre environ 8~\% de ses commandes en retard. La direction veut
maintenant comprendre ce que les clients en pensent, et ce qui se passe pendant la livraison.

\begin{enumerate}[leftmargin=*,nosep]
  \item Les retards et les incidents de livraison font-ils de mauvais avis ?
  \item Les transporteurs se valent-ils ?
  \item Où les clients sont-ils le moins satisfaits ?
  \item Peut-on repérer automatiquement, dans ce qu'ils écrivent, les clients mécontents, pour que le
        service client les traite en priorité ?
\end{enumerate}

Pour y répondre, il faut deux nouvelles sources, moins bien rangées que les tables du TP 1, les
\textbf{avis des clients}, avec un commentaire en texte libre, et les \textbf{messages JSON} que les
transporteurs envoient pour suivre chaque colis.

% =====================================================================
\section{Les données}
% =====================================================================

\begin{notion}
\textbf{Structuré, semi-structuré, non structuré.} Une donnée \emph{structurée} entre dans des
colonnes fixes (les tables du TP 1). Une donnée \emph{semi-structurée} porte sa structure avec elle,
sous forme de champs nommés, mais cette structure peut être imbriquée et varier d'un message à
l'autre, comme dans le JSON. Une donnée \emph{non structurée} n'a pas de structure exploitable
directement (texte libre, image, son). Un Lakehouse accueille les trois.
\end{notion}

\begin{center}
\small
\renewcommand{\arraystretch}{1.25}
\begin{tabularx}{\textwidth}{@{}l l X r@{}}
\toprule
\textbf{Source} & \textbf{Format} & \textbf{Contenu} & \textbf{Volume} \\
\midrule
Avis des clients & CSV & Une note de 1 à 5 et, pour 42~\% des avis, un titre et un commentaire
  saisis librement en portugais et traduits automatiquement en anglais & 100\,000 avis \\
Suivi des colis & JSON & Un message par colis, avec le transporteur, les étapes du suivi
  (prise en charge, transit, livraison) et, depuis 2018, les incidents & $\approx$ 98\,200 messages \\
Tables Silver du TP 1 & Iceberg & Commandes, articles, clients, produits, vendeurs & -- \\
\bottomrule
\end{tabularx}
\end{center}

Les avis sont les vrais avis Olist. Les clients écrivent en portugais. Pour ce TP, les titres et les commentaires ont été traduits
automatiquement en anglais par un modèle open source (OPUS-MT), en gardant la mise en forme d'origine
(fichier \cmd{data/olist/olist\_order\_reviews\_dataset\_en.csv}, à côté de l'original en portugais). La traduction est bonne dans l'ensemble, avec quelques maladresses.
Les messages de suivi sont dans le dossier \cmd{data/olist/suivi/}, un fichier par mois. Ils ont
été produits par le script \cmd{scripts/generer\_suivi.py} à partir des \textbf{vraies dates} de
remise au transporteur et de livraison des commandes. Les transporteurs, les lieux de transit et
les incidents sont simulés. Voici un message reçu en 2018 :

\begin{lstlisting}[language=json,basicstyle=\ttfamily\footnotesize]
{"colis_id": "CL0005A1A1728C", "order_id": "0005a1a1728c9d785b8e2b08b904576c",
 "emis_le": "2018-03-29 18:17:31", "version_schema": 2,
 "transporteur": {"nom": "Jadlog", "service": ".Package"},
 "evenements": [
   {"code": "PRIS_EN_CHARGE", "horodatage": "2018-03-28 00:37:42",
    "lieu": {"ville": "Sao Paulo", "etat": "SP"}},
   {"code": "EN_TRANSIT",   "horodatage": "2018-03-28 17:12:39", "lieu": {...}},
   {"code": "EN_LIVRAISON", "horodatage": "2018-03-29 13:17:31", "lieu": {...}},
   {"code": "LIVRE",        "horodatage": "2018-03-29 18:17:31", "lieu": {...}}],
 "incidents": [{"type": "COLIS_ENDOMMAGE", "horodatage": "2018-03-28 18:49:28"}]}
\end{lstlisting}

On y trouve des \textbf{objets imbriqués} (\cmd{transporteur}, \cmd{lieu}) et des \textbf{listes}
(\cmd{evenements}, \cmd{incidents}). Le flux a aussi les défauts habituels de ce type de source.
Son format a changé en janvier 2018 (les messages de 2017 n'ont ni \cmd{version\_schema}, ni
\cmd{service}, ni \cmd{incidents}), certains messages sont envoyés deux fois, d'autres arrivent
tronqués.

% =====================================================================
\section{Ce que vous allez construire}
% =====================================================================

La figure~\ref{fig:pipeline} montre le pipeline du TP 2. Comme au TP 1, les données partent des
sources, traversent les couches du Lakehouse, puis servent aux usages. Les numéros renvoient aux étapes
décrites ci-dessous.

\begin{figure}[H]
\centering
\resizebox{\textwidth}{!}{%
\begin{tikzpicture}[
  font=\sffamily\small, >=Stealth,
  barre/.style={fill=lhfonce,text=white,font=\sffamily\bfseries\small,minimum height=0.6cm,inner sep=0pt},
  cadre/.style={draw=lhfonce,line width=0.8pt},
  bloc/.style={draw=lhfonce!60,line width=0.6pt,fill=white},
  et/.style={font=\sffamily\scriptsize,text=lhfonce,align=center},
  flux/.style={->,very thick,lhfonce}]

% ===== Sources =====
\draw[cadre] (0,0) rectangle (3.2,-10.4);
\node[barre,minimum width=3.2cm,anchor=north west] at (0,0) {Sources};
\pic at (1.6,-3.6) {fichier=CSV};             \node[et] at (1.6,-4.45) {Avis des clients\\(CSV)};
\pic at (1.6,-5.6) {fichier=\{ \}};            \node[et] at (1.6,-6.45) {Suivi des colis\\(JSON)};

% ===== Plateforme Lakehouse =====
\draw[draw=lhblue,line width=2.4pt] (3.9,0) rectangle (15.9,-10.4);
\node[fill=lhblue,text=white,font=\sffamily\bfseries,minimum width=12cm,minimum height=0.7cm,anchor=north west]
  at (3.9,0) {Plateforme Lakehouse};
\draw[bloc] (4.35,-1.0) rectangle (15.2,-2.3);
\node[et,font=\sffamily\bfseries\scriptsize,anchor=west] at (4.6,-1.65) {Traitement};
\node[et,font=\sffamily\scriptsize] at (9.78,-1.65) {\logo[1.6em]{jupyter}JupyterLab\qquad\logo[1.6em]{spark}Spark\qquad\logo[1.6em]{dbt}dbt};

% Tables par couche
\pic at (5.3,-4.6)  {table=gray};      \node[et] at (5.3,-5.55)  {\textbf{landing}\\avis, suivi};
\pic at (8.1,-4.6)  {table=lhbronze};  \node[et] at (8.1,-5.55)  {\textbf{bronze}\\2 tables};
\pic at (10.9,-4.6) {table=lhsilver};  \node[et] at (10.9,-5.55) {\textbf{silver}\\+ tables du TP 1};
\pic at (13.7,-4.6) {table=lhgold};    \node[et] at (13.7,-5.55) {\textbf{gold}\\constellation};
\pic at (10.9,-6.75) {table=gray};     \node[et,anchor=west] at (11.45,-6.75) {\textbf{quarantaine}\\lignes rejetées};

% Lakehouse ouvert : catalogue et stockage
\node[fill=lhblue,text=white,font=\sffamily\bfseries\small,minimum width=10.85cm,minimum height=0.6cm]
  at (9.78,-7.75) {Lakehouse ouvert};
\draw[bloc] (4.35,-8.2) rectangle (7.85,-10.05);
\node[barre,minimum width=3.5cm,anchor=north west] at (4.35,-8.2) {Catalogue};
\node[et] at (6.1,-9.3) {\logo[1.8em]{iceberg}\\[2pt]Iceberg REST};
\draw[bloc] (8.2,-8.2) rectangle (15.2,-10.05);
\node[barre,minimum width=7.0cm,anchor=north west] at (8.2,-8.2) {Stockage ouvert (S3)};
\node[et] at (9.3,-9.3)  {\logo[1.8em]{rustfs}\\[2pt]RustFS};
\node[et] at (11.0,-9.3) {\logo[1.8em]{iceberg}\\[2pt]Iceberg};
\node[et] at (12.6,-9.3) {\logo[1.8em]{parquet}\\[2pt]Parquet};
\node[et] at (14.2,-9.3) {\logo[1.8em]{delta}\\[2pt]Delta Lake};

% ===== Usages =====
\draw[cadre] (16.6,0) rectangle (20.0,-3.2);
\node at (18.3,-1.2) {\logo[2em]{spark}}; \node[et] at (18.3,-2.15) {Requêtes SQL\\(Spark SQL)};
\draw[cadre] (16.6,-3.6) rectangle (20.0,-6.8);
\node at (18.3,-4.8) {\logo[2em]{sklearn}}; \node[et] at (18.3,-5.75) {Analyse\\du texte};
\draw[cadre] (16.6,-7.2) rectangle (20.0,-10.4);
\node at (18.3,-8.4) {\logo[2em]{streamlit}}; \node[et] at (18.3,-9.35) {Tableau\\de bord};

% ===== Infrastructure =====
\node[fill=lhfonce,text=white,minimum width=20cm,minimum height=0.75cm,anchor=north west] at (0,-10.8)
  {\logo[1.3em]{docker}\textbf{Infrastructure : conteneurs Docker}};

% ===== Flux =====
\draw[very thick,lhfonce] (3.2,-3.6) -| (4.4,-4.6) node[pastille,pos=0.15]{1};
\draw[very thick,lhfonce] (3.2,-5.6) -| (4.4,-4.6) node[pastille,pos=0.15]{2};
\draw[flux] (4.4,-4.6) -- (4.85,-4.6);
\draw[flux] (5.75,-4.6) -- (7.65,-4.6) node[pastille,midway]{3};
\draw[flux] (8.55,-4.6) -- (10.45,-4.6) node[pastille,midway]{4};
\draw[flux] (10.9,-5.95) -- (10.9,-6.2) node[pastille,right=2pt]{5};
\draw[flux] (11.35,-4.6) -- (13.25,-4.6) node[pastille,midway]{6};
\node[et,above=5pt] at (12.3,-4.6) {dbt};
\draw[very thick,lhfonce] (14.15,-4.6) -- (15.55,-4.6);
\draw[flux] (15.55,-4.6) |- (16.6,-1.6) node[pastille,pos=0.75]{7};
\draw[<->,very thick,lhfonce] (15.55,-5.2) -- (16.6,-5.2) node[pastille,midway,below=3pt]{8};
\draw[flux] (15.55,-4.6) |- (16.6,-8.8) node[pastille,pos=0.75]{9};
\end{tikzpicture}}
\caption{Pipeline du TP 2. Tous les services tournent dans des conteneurs Docker.}
\label{fig:pipeline}
\end{figure}

\subsection*{Le pipeline, étape par étape}

\begin{description}[leftmargin=1.4em,labelindent=0pt,itemsep=3pt,font=\normalfont]
\item[\etape{1}] \textbf{Dépôt des avis} (\nb{01\_bronze\_avis\_suivi}). Le fichier CSV des avis est copié tel quel dans
  \cmd{s3://landing/olist/avis/}.
\item[\etape{2}] \textbf{Dépôt du suivi} (\nb{01\_bronze\_avis\_suivi}). Les messages des transporteurs, un fichier par mois et un
  message JSON par ligne, sont copiés dans la zone d'atterrissage (\cmd{suivi/}).
\item[\etape{3}] \textbf{Ingestion en Bronze} (\nb{01\_bronze\_avis\_suivi}). Spark lit les avis en texte, en autorisant les
  commentaires sur plusieurs lignes, et garde chaque message JSON \emph{comme un texte}, sans le
  découper. Les deux tables reçoivent les colonnes \cmd{\_source} et \cmd{\_ingere\_le}.
\item[\etape{4}] \textbf{Nettoyage} (\nb{02\_silver\_texte\_json}). Les avis sont typés et leur
  commentaire est normalisé (minuscules, sans accents ni ponctuation). Les messages JSON sont lus avec
  un schéma écrit à l'avance, dédoublonnés, puis \og dépliés \fg{}, avec une ligne par colis et une
  ligne par étape du suivi.
\item[\etape{5}] \textbf{Quarantaine} (\nb{02\_silver\_texte\_json}). Les avis dont la commande est inconnue et les messages
  tronqués sont isolés.
\item[\etape{6}] \textbf{Modélisation avec dbt} (\nb{03\_gold\_dbt\_etoile}). Un projet dbt construit le schéma en constellation à partir
  des tables Silver des TP 1 et 2, puis vérifie chaque table par des tests de qualité.
\item[\etape{7}] \textbf{Analyse} (\nb{03\_gold\_dbt\_etoile}). Des requêtes SQL sur ce schéma répondent aux questions de la direction.
\item[\etape{8}] \textbf{Analyse du texte} (\nb{04\_nlp\_avis}). Un modèle apprend à reconnaître un avis négatif à partir
  de son commentaire, puis écrit ses prédictions dans \cmd{gold.avis\_sentiment}.
\item[\etape{9}] \textbf{Restitution} (\nb{dashboard\_tp2.py}). Le tableau de bord présente la satisfaction des clients et
  permet de tester le modèle sur un nouveau commentaire.
\end{description}

Il s'agit toujours d'un pipeline \textbf{ELT}. Les étapes \etape{1} à \etape{3} extraient et chargent
sans rien transformer, et toutes les transformations (\etape{4} à \etape{8}) ont lieu dans le Lakehouse.

\subsection*{La couche Gold en schéma en constellation}

Au TP 1, la couche Gold était une table large, \cmd{gold.commandes}. Elle est remplacée ici par une
\textbf{modélisation dimensionnelle} (figure~\ref{fig:constellation}), qui distingue deux sortes de tables.

\begin{itemize}[leftmargin=*,itemsep=1pt]
\item Une table de \textbf{faits} enregistre des événements mesurables (une vente, un avis, une
  livraison). Son \textbf{grain} dit ce que représente une ligne.
\item Une table de \textbf{dimensions} décrit le contexte qui sert à filtrer et à regrouper (quand,
  qui, quoi). Chaque fait la désigne par une \textbf{clé}.
\end{itemize}

Une table de faits entourée de ses dimensions forme un \textbf{schéma en étoile}. Ici, il y a
\emph{trois} tables de faits, qui partagent les dimensions \cmd{dim\_date} et \cmd{dim\_client} (des
\emph{dimensions conformes}). Trois étoiles reliées entre elles forment un \textbf{schéma en
constellation} (\emph{fact constellation}, ou \emph{galaxy schema}). Une même question (\og par état
du client \fg{}, \og par mois \fg{}) se pose ainsi de la même façon aux ventes, aux avis et aux
livraisons. Le projet dbt s'appelle \cmd{dbt\_etoile} car chaque fait y est modélisé en étoile.

\begin{figure}[H]
\centering
\begin{tikzpicture}[
  font=\sffamily\small, >=Stealth,
  fait/.style={draw=lhgold,thick,rounded corners=3pt,fill=lhgold!12,text width=3.6cm,align=left,
               inner sep=5pt,font=\sffamily\footnotesize},
  dim/.style={draw=lhblue,thick,rounded corners=3pt,fill=lhblue!6,text width=3.0cm,align=left,
              inner sep=5pt,font=\sffamily\footnotesize},
  lien/.style={thick,lhblue!70}]
\node[fait] (fv) at (0,0)    {\textbf{fait\_ventes}\\ \textit{grain : un article}\\[2pt] \texttt{prix}, \texttt{frais\_port}\\ \texttt{statut}, \texttt{en\_retard}};
\node[fait] (fa) at (5.0,0)  {\textbf{fait\_avis}\\ \textit{grain : un avis}\\[2pt] \texttt{note}, \texttt{est\_negatif}\\ \texttt{nb\_mots}, \texttt{en\_retard}};
\node[fait] (fl) at (10.0,0) {\textbf{fait\_livraisons}\\ \textit{grain : un colis}\\[2pt] \texttt{transporteur}\\ \texttt{nb\_incidents}, \texttt{duree}};

\node[dim] (dp) at (-1.3,3.1) {\textbf{dim\_produit}\\ \underline{\texttt{product\_id}}\\ \texttt{categorie}};
\node[dim] (dv) at (2.4,3.1)  {\textbf{dim\_vendeur}\\ \underline{\texttt{seller\_id}}\\ \texttt{ville}, \texttt{etat}};
\node[dim] (dc) at (8.4,3.1)  {\textbf{dim\_client}\\ \underline{\texttt{customer\_id}}\\ \texttt{ville}, \texttt{etat}};
\node[dim] (dd) at (5.0,-3.1) {\textbf{dim\_date}\\ \underline{\texttt{date\_id}}\\ \texttt{annee\_mois}, \dots};

\draw[lien] (dp.south) -- ([xshift=-8mm]fv.north);
\draw[lien] (dv.south) -- ([xshift=8mm]fv.north);
\draw[lien] (dc.south) -- (fv.north east);
\draw[lien] (dc.south) -- (fa.north);
\draw[lien] (dc.south) -- (fl.north);
\draw[lien] (dd.north) -- (fv.south);
\draw[lien] (dd.north) -- (fa.south);
\draw[lien] (dd.north) -- (fl.south);
\end{tikzpicture}
\caption{Le schéma en constellation de la couche Gold. Trois tables de faits (en jaune) reliées aux
dimensions (en bleu) par leurs clés (soulignées). Chaque fait garde aussi le numéro de commande
(\texttt{order\_id}), qui permet de rapprocher deux faits, avec prudence, car leurs grains sont différents.}
\label{fig:constellation}
\end{figure}

\begin{center}
\small
\renewcommand{\arraystretch}{1.3}
\begin{tabularx}{\textwidth}{@{}>{\bfseries}p{2.2cm} X X@{}}
\toprule
\normalfont\textbf{Couche} & \textbf{Rôle} & \textbf{Règle d'or} \\
\midrule
Bronze & Copie des sources, JSON compris, avec \cmd{\_source} et \cmd{\_ingere\_le} &
  On ne découpe pas, on ne corrige pas, et un message tronqué est gardé tel quel. \\
Silver & Données typées, dédoublonnées, dépliées, texte normalisé &
  Le schéma est écrit, jamais deviné. Une ligne invalide va en quarantaine. \\
Gold & Schéma en constellation, construit et testé par dbt &
  Chaque table a un grain clair et passe ses tests avant d'être utilisée. \\
\bottomrule
\end{tabularx}
\end{center}

% =====================================================================
\section{Déroulé}
% =====================================================================

\textbf{Avant de commencer, récupérez les énoncés et copiez-les dans votre dossier personnel.} Dans
JupyterLab, ouvrez un terminal (\emph{File > New > Terminal}) et tapez :
\begin{lstlisting}[style=shell]
git pull
cp -r notebooks/tp2 notebooks/mon_travail/
cp apps/dashboard_tp2.py notebooks/mon_travail/tp2/
\end{lstlisting}
Le dossier \nb{notebooks/mon\_travail/tp2/} contient alors les quatre notebooks, le projet dbt
\nb{dbt\_etoile/} et le tableau de bord. Travaillez uniquement dans ce dossier, et exécutez les
notebooks \textbf{dans l'ordre}. Les cellules marquées \cmd{\# À COMPLÉTER} sont à écrire. Dans le
projet dbt, ce sont les fichiers qui portent la mention \cmd{-{}- À COMPLÉTER}.

\begin{attention}
\textbf{Un seul processus Spark à la fois.} Un notebook ouvert, une commande \cmd{dbt} et le tableau
de bord démarrent chacun leur propre session Spark. Avant de passer de l'un à l'autre, arrêtez les
noyaux des notebooks (panneau \emph{Running Terminals and Kernels} > \emph{Shut Down All}).
\end{attention}

L'aide-mémoire de la dernière page rassemble les fonctions Spark et les commandes dbt de ce TP.

% =====================================================================
\section{Partie par partie}
% =====================================================================

\subsection*{Partie 1 -- Bronze \hfill \nb{01\_bronze\_avis\_suivi}}

Vous déposez les deux sources dans la zone d'atterrissage, vous mesurez le piège des retours à la
ligne dans les commentaires, puis vous écrivez les tables \cmd{bronze.avis} et \cmd{bronze.suivi\_brut}.

\begin{notion}
\textbf{Un CSV n'a pas toujours un enregistrement par ligne.} Quand un client écrit son commentaire
sur plusieurs lignes, le retour à la ligne se trouve \emph{entre guillemets}, à l'intérieur du champ.
Une lecture naïve le prend pour une fin d'enregistrement et coupe l'avis en morceaux. L'option
\cmd{multiLine} de Spark respecte les guillemets.
\end{notion}

\begin{notion}
\textbf{Garder le JSON brut.} En Bronze, chaque message est stocké comme un simple texte. Si le
transporteur change de format, l'ingestion continue de fonctionner, et un message mal formé n'est
jamais perdu. Le découpage se fait en Silver.
\end{notion}

\begin{resultat}
\cmd{bronze.avis} contient exactement 100\,000 avis, avec des notes de 1 à 5 uniquement (la lecture
naïve en trouvait davantage). \cmd{bronze.suivi\_brut} contient environ 98\,200 messages.
\end{resultat}

\capture{tp2_01_notebook}{Notebook \texttt{01\_bronze\_avis\_suivi} : lecture naïve et lecture multi-lignes des avis.}{fig:nb01}
\capture{tp2_01_resultat}{Résultat attendu de la partie 1 : le bilan des deux tables Bronze.}{fig:res01}

\subsection*{Partie 2 -- Silver \hfill \nb{02\_silver\_texte\_json}}

Vous trouvez la bonne clé des avis, vous normalisez leur texte, puis vous découpez les messages JSON
en deux tables, avec une ligne par colis dans la première et une ligne par étape du suivi dans la seconde.

\begin{notion}
\textbf{Lire le JSON avec un schéma écrit.} \cmd{F.from\_json(colonne, schema)} transforme le texte
en structure. Le schéma couvre les deux versions du format, et un champ absent vaut \cmd{null}. On
accède ensuite aux champs imbriqués par un point (\cmd{transporteur.nom}) et aux éléments d'une liste
par leur position (\cmd{incidents[0]}).
\end{notion}

\begin{notion}
\textbf{Déplier une liste.} \cmd{F.explode} crée une ligne par élément d'une liste
(figure~\ref{fig:explode}). C'est ainsi qu'un message imbriqué devient une table.
\end{notion}

\begin{figure}[H]
\centering
\begin{tikzpicture}[font=\sffamily\footnotesize, >=Stealth,
  b/.style={draw=#1,thick,rounded corners=3pt,fill=#1!7,align=left,inner sep=5pt}]
\node[b=lhbronze,text width=4.6cm] (m) at (0,0) {\textbf{un message} (une ligne)\\[3pt]
  \texttt{colis\_id} : \texttt{CL0005A1A1728C}\\[1pt]
  \texttt{evenements} : [\\
  \quad \texttt{PRIS\_EN\_CHARGE},\\ \quad \texttt{EN\_TRANSIT},\\ \quad \texttt{EN\_LIVRAISON},\\ \quad \texttt{LIVRE} ]};
\node[b=lhsilver,text width=6.4cm] (t) at (8.2,0) {\textbf{suivi\_evenements} (quatre lignes)\\[3pt]
  \scriptsize\ttfamily
  \begin{tabular}{@{}lll@{}}
  CL0005A1A1728C & PRIS\_EN\_CHARGE & 28/03 \\
  CL0005A1A1728C & EN\_TRANSIT     & 28/03 \\
  CL0005A1A1728C & EN\_LIVRAISON   & 29/03 \\
  CL0005A1A1728C & LIVRE           & 29/03 \\
  \end{tabular}};
\draw[->,thick] (m) -- node[above]{\texttt{F.explode}} (t);
\end{tikzpicture}
\caption{\texttt{F.explode}. Une liste de quatre événements devient quatre lignes, et les colonnes du
message (\texttt{colis\_id}) sont répétées sur chaque ligne.}
\label{fig:explode}
\end{figure}

\begin{attention}
Sur une liste \cmd{null} (un message de 2017, sans \cmd{incidents}), \cmd{F.size} renvoie \cmd{-1},
et non \cmd{0}. Ce genre de détail fausse une moyenne sans provoquer d'erreur. Il faut donc le traiter
explicitement, et le tester (partie 3).
\end{attention}

\begin{resultat}
Le bilan explique chaque écart entre Bronze et Silver, avec quelques avis orphelins (leur commande est
en quarantaine depuis le TP 1), environ 200 messages tronqués et environ 600 messages reçus deux
fois. \cmd{silver.suivi\_colis} compte environ 97\,450 colis, \cmd{silver.suivi\_evenements} environ
450\,000 étapes.
\end{resultat}

\capture{tp2_02_notebook}{Notebook \texttt{02\_silver\_texte\_json} : lecture du JSON avec un schéma explicite.}{fig:nb02}
\capture{tp2_02_resultat}{Résultat attendu de la partie 2 : le bilan Bronze / Silver.}{fig:res02}

\subsection*{Partie 3 -- Gold avec dbt \hfill \nb{03\_gold\_dbt\_etoile}}

Vous complétez le projet dbt (un modèle de dimension, deux modèles de faits et leurs tests), vous
construisez toute la couche Gold avec \cmd{dbt build}, puis vous interrogez ces tables en SQL.

\begin{notion}
\textbf{dbt (\emph{data build tool}).} Chaque table Gold est décrite par un fichier SQL qui contient
une simple requête \cmd{SELECT}, appelé \emph{modèle}. dbt en fait une table, dans le bon ordre, car il
déduit les dépendances des expressions \cmd{\{\{ source(...) \}\}} (une table Silver) et
\cmd{\{\{ ref(...) \}\}} (un autre modèle). Les \textbf{tests} (\cmd{unique}, \cmd{not\_null},
\cmd{accepted\_values}, \cmd{relationships}, ou une requête SQL qui renvoie les lignes en erreur)
sont lancés juste après la construction de chaque table. Si l'un échoue, les tables qui en
dépendent ne sont pas construites. Le projet, versionné comme du code, documente lui-même la couche
Gold.
\end{notion}

\needspace{22\baselineskip}
\paragraph{Le projet \cmd{dbt\_etoile}.} Il se trouve dans \cmd{notebooks/mon\_travail/tp2/dbt\_etoile/}.
Les fichiers marqués \og à compléter \fg{} contiennent des trous :

\begin{lstlisting}[style=shell]
dbt_etoile/
|-- dbt_project.yml       # configuration : tables Iceberg dans lakehouse.gold
|-- models/
|   |-- sources.yml       # tables Silver lues par le projet
|   |-- schema.yml        # colonnes et tests                  <- à compléter
|   |-- dimensions/
|   |   |-- dim_date.sql     dim_client.sql     dim_vendeur.sql
|   |   `-- dim_produit.sql                                    <- à compléter
|   `-- faits/
|       |-- fait_ventes.sql
|       |-- fait_avis.sql                                      <- à compléter
|       `-- fait_livraisons.sql                                <- à compléter
|-- tests/                # tests SQL : renvoient les lignes en erreur
|   |-- fait_ventes_grain_unique.sql
|   `-- nb_incidents_positif.sql
`-- target/               # créé par dbt (SQL compilé, documentation)
\end{lstlisting}

La connexion à Spark est déjà configurée dans l'image (\cmd{profiles.yml}), et \cmd{dbt debug
-{}-connection} doit répondre \cmd{All checks passed!}. Les commandes se lancent depuis le notebook
(préfixe \cmd{!}) ou depuis un terminal JupyterLab :

\begin{lstlisting}[style=shell]
cd notebooks/mon_travail/tp2
dbt build --project-dir dbt_etoile            # tous les modèles et tous les tests
dbt run   --project-dir dbt_etoile --select fait_avis   # un seul modèle
\end{lstlisting}

\begin{resultat}
\cmd{dbt build} se termine par \cmd{Completed successfully} et \cmd{PASS=31}, soit 7 modèles et 24 tests
(22 déclarés dans \cmd{schema.yml}, 2 écrits en SQL dans \cmd{tests/}). Environ 55~\% des avis sont
négatifs quand la livraison est en retard, contre environ 10~\% sinon.
\end{resultat}

\capture{tp2_03_dbt_build}{Résultat attendu de \texttt{dbt build} : 7 modèles et 24 tests, sans erreur (\texttt{PASS=31}).}{fig:dbtbuild}
\paragraph{Consulter le graphe des dépendances (\emph{lineage}).} dbt génère un site de
documentation, avec la description de chaque table et de chaque colonne, le SQL, les tests et le graphe
des dépendances.
\begin{enumerate}[leftmargin=*,nosep]
  \item Lancez la cellule \cmd{!dbt docs generate -{}-static -{}-project-dir dbt\_etoile} du notebook.
        Elle crée un fichier unique, \cmd{dbt\_etoile/target/static\_index.html}.
  \item Dans l'explorateur de fichiers de JupyterLab, ouvrez \cmd{dbt\_etoile/target/}, faites un
        clic droit sur \cmd{static\_index.html} puis \emph{Open in New Browser Tab}.
  \item Dans la page qui s'ouvre, cliquez sur un modèle (par exemple \cmd{fait\_avis}) dans le menu
        de gauche, puis sur l'icône bleue \emph{View Lineage Graph}, en bas à droite.
  \item Le graphe montre les tables Silver (sources, en vert) à gauche et les modèles Gold à droite.
        Le champ \emph{-{}-select} en bas de la fenêtre filtre le graphe, par exemple
        \cmd{+fait\_avis} pour ne voir que ce dont \cmd{fait\_avis} dépend.
\end{enumerate}
Le graphe s'arrête aux tables Silver, car les couches Bronze et Silver sont construites par Spark, pas
par dbt, qui ne les connaît que comme des \emph{sources} déclarées dans \cmd{sources.yml}. Relancez
\cmd{dbt docs generate} après chaque modification du projet, puis actualisez la page.

\capture{tp2_03_lineage}{La documentation générée par dbt : le graphe des dépendances (\emph{lineage}) de la couche Gold.}{fig:lineage}
\capture{tp2_03_resultat}{Résultat attendu de la partie 3 : avis négatifs selon le retard de livraison (question Q1).}{fig:res03}

\begin{attention}
Une corrélation n'est pas une cause. Les colis qui ont subi un incident reçoivent de moins bonnes
notes, mais ils sont aussi plus souvent en retard. Le notebook vous fait refaire la comparaison
\emph{à retard égal} pour démasquer cette \textbf{variable de confusion}.
\end{attention}

\begin{attention}
Une jointure entre deux faits de grains différents multiplie les lignes. Une commande de trois
articles a trois lignes dans \cmd{fait\_ventes}, et son avis est compté trois fois. Le notebook vous
fait trouver et corriger ce piège.
\end{attention}

\subsection*{Partie 4 -- Analyse du texte \hfill \nb{04\_nlp\_avis}}

Vous entraînez un modèle qui estime, à partir du commentaire, la probabilité qu'un avis soit
négatif (note 1 ou 2). Vous l'interprétez, vous étudiez ses erreurs, puis vous écrivez ses
prédictions dans \cmd{gold.avis\_sentiment}.

\begin{notion}
\textbf{TF-IDF.} Pour qu'un modèle lise du texte, chaque commentaire devient un vecteur, avec un poids
par terme du vocabulaire. Ce poids est élevé quand le terme est fréquent \emph{dans ce commentaire}
et rare \emph{dans l'ensemble des commentaires}. Les termes peuvent être des mots seuls ou des groupes
de deux mots (\emph{n-grammes}), car \og not recommend \fg{} ne dit pas la même chose que
\og recommend \fg{}. Une régression logistique apprend ensuite un coefficient par terme, et le modèle
reste lisible.
\end{notion}

\begin{figure}[H]
\centering
\resizebox{\textwidth}{!}{%
\begin{tikzpicture}[font=\sffamily\footnotesize, >=Stealth, node distance=0.9cm,
  b/.style={draw=#1,thick,rounded corners=3pt,fill=#1!7,align=center,inner sep=5pt,text width=3.0cm,minimum height=1.5cm}]
\node[b=gray] (a) {\og I do NOT recommend!!\\ Arrived broken. \fg{}};
\node[b=lhsilver,right=of a] (b) {texte normalisé\\[2pt] \texttt{\scriptsize i do not recommend arrived broken}};
\node[b=lhblue,right=of b] (c) {vecteur TF-IDF\\[2pt] \scriptsize not recommend : 0,52\\ broken : 0,48\\ \dots};
\node[b=lhblue,right=of c] (d) {régression\\ logistique};
\node[b=lhgold,right=of d] (e) {probabilité d'un\\ avis négatif :\\ \textbf{0,98}};
\draw[->,thick] (a) -- (b);
\draw[->,thick] (b) -- (c);
\draw[->,thick] (c) -- (d);
\draw[->,thick] (d) -- (e);
\end{tikzpicture}}
\caption{Du commentaire à la prédiction, avec la normalisation (partie 2), vectorisation TF-IDF, puis
classification. Les valeurs sont données à titre d'illustration.}
\label{fig:tfidf}
\end{figure}

\begin{attention}
Les listes de \emph{mots vides} (articles, pronoms, prépositions\dots) toutes faites contiennent
souvent la négation. Celle de scikit-learn retire \og not \fg{}, \og no \fg{} et \og never \fg{}. Avec
elle, \og i do not recommend \fg{} devient \og recommend \fg{} et passe pour un avis positif, ce que le
notebook vous fait constater.
\end{attention}

\begin{resultat}
Une AUC ROC d'environ 0,96 et un rappel d'environ 90~\% sur les avis négatifs. Le client écrit
lui-même ce qui ne va pas, et la tâche est bien plus facile que la prévision des retards du TP 1.
\end{resultat}

\capture{tp2_04_notebook}{Notebook \texttt{04\_nlp\_avis} : entraînement et évaluation du modèle TF-IDF.}{fig:nb04}
\capture{tp2_04_resultat}{Résultat attendu de la partie 4 : les termes qui pèsent le plus vers un avis négatif ou positif.}{fig:res04}

% =====================================================================
\section{Le tableau de bord}\label{sec:dashboard}
% =====================================================================

Ouvrez \nb{notebooks/mon\_travail/tp2/dashboard\_tp2.py} dans JupyterLab. Les requêtes sur la couche Gold
et le filtre de période sont fournis. Complétez les cinq éléments demandés :

\begin{enumerate}[leftmargin=*,nosep]
  \item une ligne de quatre indicateurs : note moyenne, part d'avis négatifs, part d'avis commentés,
        part des colis avec incident ;
  \item la part d'avis négatifs, mois par mois ;
  \item la part d'avis négatifs selon que la livraison était à l'heure ou en retard ;
  \item la part de colis avec incident, par transporteur ;
  \item une zone où l'on saisit un commentaire et où le modèle de la partie 4 affiche sa probabilité
        d'être négatif, suivie des avis à traiter en priorité.
\end{enumerate}

Arrêtez les noyaux des notebooks, puis lancez le tableau de bord depuis un terminal :
\begin{lstlisting}[style=shell]
streamlit run notebooks/mon_travail/tp2/dashboard_tp2.py --server.address 0.0.0.0
\end{lstlisting}
Ouvrez ensuite \url{http://localhost:8501} dans votre navigateur.

\begin{astuce}
Inutile de relancer Streamlit à chaque modification. Enregistrez le fichier dans JupyterLab
(\cmd{Ctrl+S}), et un bandeau \emph{Source file changed} apparaît en haut à droite de la page. Cliquez
sur \emph{Always rerun}, la page se met alors à jour à chaque enregistrement. De même, chaque
action dans la page (changer la période, saisir un commentaire) réexécute tout le script de haut en
bas, c'est le fonctionnement normal de Streamlit. Les résultats des requêtes sont gardés en cache
cinq minutes (\cmd{@st.cache\_data}). Pour forcer leur relecture, ouvrez le menu \emph{\(\vdots\)} >
\emph{Clear cache}, puis tapez \cmd{R}. Pour arrêter le tableau de bord, faites \cmd{Ctrl+C} dans le terminal.
\end{astuce}

\capture{tp2_dashboard_haut}{Tableau de bord attendu : les indicateurs, l'évolution mensuelle et l'effet des retards.}{fig:dash1}
\capture{tp2_dashboard_bas}{Tableau de bord attendu : les transporteurs et le test d'un commentaire par le modèle.}{fig:dash2}

% =====================================================================
\section{Livrables}\label{sec:livrables}
% =====================================================================

\subsection*{Le compte rendu}

Au format PDF, nommé \cmd{NOM\_Prenom\_TP2.pdf}, il ne recopie pas le code mais explique vos choix
et vos résultats. Il contient :
\begin{enumerate}[leftmargin=*,itemsep=2pt]
  \item \textbf{Architecture} : le schéma de votre pipeline, et ce qui change par rapport au TP 1
        quand les données ne sont pas structurées.
  \item \textbf{Ingestion et qualité} : le piège du CSV multi-lignes, le choix de garder le JSON brut,
        la clé des avis, les règles de quarantaine et le bilan Bronze / Silver.
  \item \textbf{Schéma en constellation} : le grain de chaque fait, le rôle des dimensions, une capture du
        \emph{lineage} dbt, les tests que vous avez écrits et ce qu'ils protègent.
  \item \textbf{Réponses à la direction} : les résultats des requêtes SQL, interprétés.
  \item \textbf{Analyse du texte} : la méthode, les métriques, les termes qui pèsent le plus, l'effet
        des mots vides et l'analyse des erreurs.
  \item \textbf{Tableau de bord} : une capture d'écran et deux ou trois enseignements.
  \item \textbf{Réponses aux questions} des notebooks, rédigées et argumentées.
  \item \textbf{Limites et améliorations} : ce que vous feriez avec plus de temps, ou en production
        (par exemple des modèles de langue pré-entraînés, ou le traitement des avis mixtes).
\end{enumerate}

\clearpage
% =====================================================================
\section*{Annexe -- Aide-mémoire}\label{sec:aidememoire}
% =====================================================================

Les imports sont \cmd{from pyspark.sql import functions as F, Window}. Les opérations du TP 1 restent valables.

\begin{center}
\small
\renewcommand{\arraystretch}{1.22}
\begin{tabularx}{\textwidth}{@{}>{\raggedright\arraybackslash}p{4.1cm} >{\raggedright\arraybackslash\ttfamily\footnotesize}X@{}}
\toprule
\normalfont\textbf{Je veux\dots} & \normalfont\textbf{Code} \\
\midrule
\multicolumn{2}{@{}l}{\textbf{\textcolor{lhblue}{Lire}}} \\
CSV avec du texte sur plusieurs lignes & spark.read.option("header", True).option("multiLine", True).option("escape", '"').csv(chemin) \\
Fichiers ligne par ligne, en texte & spark.read.text(chemin) \quad \textrm{(colonne }value\textrm{)} \\
Laisser Spark deviner un JSON & spark.read.json(chemin).printSchema() \\
\midrule
\multicolumn{2}{@{}l}{\textbf{\textcolor{lhblue}{JSON et listes}}} \\
Lire un JSON avec un schéma & F.from\_json("value", "colis\_id STRING, evenements ARRAY<STRUCT<code: STRING>>") \\
Champ imbriqué & F.col("transporteur.nom") \quad F.col("m.*") \\
Élément d'une liste & F.col("incidents")[0]["type"] \\
Nombre d'éléments & F.size("evenements") \quad \textrm{(-1 si la liste est nulle)} \\
Une ligne par élément & df.select("colis\_id", F.explode("evenements").alias("e")) \\
\midrule
\multicolumn{2}{@{}l}{\textbf{\textcolor{lhblue}{Texte}}} \\
Minuscules, sans accents & F.lower(c) \quad F.translate(c, "áàãéêíóôõúç", "aaaeeiooouc") \\
Remplacer par une expression régulière & F.regexp\_replace(c, r"[\^{}a-z\textbackslash s]", " ") \\
Coller deux textes & F.concat\_ws(" ", "titre", "message") \\
Découper en mots & F.split(c, " ") \\
\midrule
\multicolumn{2}{@{}l}{\textbf{\textcolor{lhblue}{dbt (dans un terminal, depuis \texttt{mon\_travail/tp2})}}} \\
Vérifier la connexion & dbt debug -{}-connection -{}-project-dir dbt\_etoile \\
Construire un modèle & dbt run -{}-project-dir dbt\_etoile -{}-select dim\_produit \\
Tout construire et tester & dbt build -{}-project-dir dbt\_etoile \\
Générer la documentation & dbt docs generate -{}-static -{}-project-dir dbt\_etoile \\
Désigner une table Silver / un modèle & \{\{ source('silver', 'orders') \}\} \quad \{\{ ref('dim\_date') \}\} \\
\midrule
\multicolumn{2}{@{}l}{\textbf{\textcolor{lhblue}{scikit-learn}}} \\
Vectoriser du texte & TfidfVectorizer(stop\_words=MOTS\_VIDES, ngram\_range=(1, 2), min\_df=5) \\
Enchaîner les étapes & Pipeline([("tfidf", ...), ("modele", LogisticRegression(...))]) \\
Probabilité de la classe 1 & pipe.predict\_proba(textes)[:, 1] \\
Sauvegarder un modèle & joblib.dump(pipe, "modele.joblib") \quad joblib.load(...) \\
\bottomrule
\end{tabularx}
\end{center}

\end{document}
